\documentclass[12pt, a4paper]{exam}
% Corrigé rédigé pour le site Ressources NSI et SNT (CC BY-SA 4.0).
% Compilation : pdflatex (deux passes).
\usepackage[T1]{fontenc}
\usepackage[utf8]{inputenc}
\usepackage{lmodern}
\usepackage[french]{babel}
\usepackage[a4paper,top=2.3cm,bottom=2.8cm,left=2cm,right=2cm,headsep=0.6cm,footskip=1.2cm]{geometry}
\usepackage{xcolor}
\usepackage{listings}
\usepackage{graphicx}
\usepackage{array}
\usepackage{amsmath}
\usepackage{amssymb}
\usepackage{enumitem}
\usepackage{needspace}
\usepackage{tikz}
\usetikzlibrary{positioning,arrows.meta,calc}
\usepackage[most]{tcolorbox}

% Coloration du code : même style que les sujets
\definecolor{codegreen}{rgb}{0,0.6,0}
\definecolor{codegray}{rgb}{0.5,0.5,0.5}
\definecolor{codepurple}{rgb}{0.58,0,0.82}
\definecolor{backcolour}{rgb}{0.95,0.95,0.92}
\lstdefinestyle{mystyle}{
    backgroundcolor=\color{backcolour},
    commentstyle=\color{codegreen}\itshape,
    keywordstyle=\color{magenta}\bfseries,
    numberstyle=\tiny\color{codegray},
    stringstyle=\color{codepurple},
    basicstyle=\footnotesize\ttfamily,
    breakatwhitespace=false,
    breaklines=true,
    postbreak=\mbox{\textcolor{codegray}{$\hookrightarrow$}\space},
    captionpos=b,
    keepspaces=true,
    columns=fullflexible,
    numbers=left,
    numbersep=5pt,
    xleftmargin=1.5em,
    showspaces=false,
    showstringspaces=false,
    showtabs=false,
    tabsize=4,
    upquote=true,
    extendedchars=true,
    literate={à}{{\`a}}1 {â}{{\^a}}1 {é}{{\'e}}1 {è}{{\`e}}1 {ê}{{\^e}}1 {ë}{{\"e}}1
             {î}{{\^i}}1 {ï}{{\"i}}1 {ô}{{\^o}}1 {ù}{{\`u}}1 {û}{{\^u}}1 {ç}{{\c c}}1
             {É}{{\'E}}1 {À}{{\`A}}1 {œ}{{\oe}}1 {…}{{\ldots}}1
}
\lstset{style=mystyle, language=Python}
\newcommand{\code}[1]{\texttt{#1}}

% Encadrés du corrigé (lisibles en noir et blanc)
\tcbset{corrbase/.style={enhanced, breakable, sharp corners, left=2mm, right=2mm, top=1mm, bottom=1mm,
  fonttitle=\bfseries\small, coltitle=black, before skip=2mm, after skip=3mm}}
\newtcolorbox{reponse}[1][Réponse]{corrbase, unbreakable, colback=gray!6, colframe=black, colbacktitle=gray!22, boxrule=0.8pt, title={#1}}
\newtcolorbox{erreur}[1][Erreur fréquente]{corrbase, colback=white, colframe=black, boxrule=0.4pt, leftrule=3pt, title={#1}, colbacktitle=white, attach title to upper={\ :\ }}
\newtcolorbox{retenir}[1][À retenir]{corrbase, colback=white, colframe=black, boxrule=0.4pt, borderline={0.4pt}{1.2pt}{black}, title={#1}, colbacktitle=white, attach title to upper={\ :\ }}

% En-tête d'une question : numéro, points, rappel de la question
\newcommand{\Q}[3]{\par\needspace{5\baselineskip}\bigskip\noindent{\large\bfseries Question #1}\ifx&#2&\else\ \textnormal{(#2)}\fi\par\nopagebreak\smallskip\noindent\textit{#3}\par\nopagebreak\smallskip}
\newcommand{\Qp}[2]{\par\medskip\noindent\textbf{#1}\ \textit{#2}\par\nopagebreak\smallskip}
\newcommand{\methode}{\par\noindent\textbf{Méthode.}\ }
\newcommand{\application}{\par\noindent\textbf{Application.}\ }

% Arbres
\tikzset{arbre/.style={every node/.style={draw,circle,minimum size=7mm,inner sep=1pt,font=\small},level distance=11mm, thick},
  marque/.style={fill=gray!35, line width=1.6pt},
  titre/.style={draw=none, rectangle, inner sep=2pt}}

\raggedbottom
\setlength{\parindent}{0pt}
\setlength{\parskip}{3pt}
\setlist{itemsep=1pt, topsep=2pt}
\newtcolorbox{remarque}[1][Remarque sur l'énoncé]{corrbase, colback=white, frame hidden, borderline={0.6pt}{0pt}{black, dashed}, title={#1}, colbacktitle=white, attach title to upper={\ :\ }}

\title{Corrigé --- CC Term NSI - Arbre et graphes}
\author{Babash - NSI Tani Malandi et Johanne Bakalara}
\date{Décembre 2025 \textperiodcentered{} Durée du contrôle : 2 h \textperiodcentered{} 19 points}

\newcommand{\arbreUn}[1][]{%
\begin{tikzpicture}[arbre, level 1/.style={sibling distance=26mm}, level 2/.style={sibling distance=14mm}]
\node {13} child {node {6} child {node {7}} child {node {5} child[missing] child {node {1}}}} child {node {4}};
#1
\end{tikzpicture}}
\newcommand{\arbreDeux}[1][]{%
\begin{tikzpicture}[arbre, level 1/.style={sibling distance=26mm}, level 2/.style={sibling distance=13mm}]
\node {6} child {node {4} child {node {3}} child {node {5}}} child {node {8} child {node {7}} child {node {9}}};
#1
\end{tikzpicture}}

\begin{document}
\pagestyle{foot}
\footer{}{\footnotesize Babash - NSI Tani Malandi et Johanne Bakalara \textperiodcentered{} CC BY-SA 4.0 \textperiodcentered{} page \thepage}{}
\maketitle
\thispagestyle{foot}
\begin{center}\footnotesize\itshape Exercice sur les arbres d'après l'exercice 1 du sujet de NSI du baccalauréat 2023 (jour 2, Polynésie, code 23-NSIJ2PO1).\end{center}

\begin{retenir}[Comment lire ce corrigé]
Pour chaque question : la question rappelée en italique, la réponse attendue dans un cadre gris, puis l'explication (la méthode en étapes, puis son application). Les cadres « Erreur fréquente » montrent les pièges classiques. Le sujet demande des réponses rédigées et justifiées : les justifications données ici sont celles qu'on attend sur ta copie.
\end{retenir}

\begin{center}
\arbreUn[{\node[titre, font=\bfseries] at (0,0.8) {Arbre 1};}]
\hspace{1.5cm}
\arbreDeux[{\node[titre, font=\bfseries] at (0,0.8) {Arbre 2};}]
\end{center}

\begin{retenir}[Les conventions de ce corrigé]
La \textbf{profondeur} d'un nœud est le nombre d'arêtes (de traits) entre la racine et lui : la racine a la profondeur 0. La \textbf{hauteur} d'un arbre est la plus grande profondeur d'un de ses nœuds. Certains cours comptent les nœuds au lieu des arêtes : il faut alors ajouter 1 aux hauteurs données ici. Au bac, la convention est rappelée dans l'énoncé.
\end{retenir}

\section*{1. Notions fondamentales sur les arbres}

\textbf{Question 1 : question de cours} (4 points au total).

\Q{1 (a)}{1 point}{Recopier et compléter le texte avec : au maximum, au minimum, exactement, feuille, racine, sous-arbre gauche, sous-arbre droit.}
\begin{reponse}
\begin{enumerate}[label=(\alph*)]
\item Le nœud initial est appelé \textbf{racine}.
\item Un nœud qui n'a pas d'enfants est appelé \textbf{feuille}.
\item Un arbre binaire est un arbre dans lequel chaque nœud a \textbf{au maximum} deux enfants.
\item Un arbre binaire de recherche est un arbre binaire dans lequel tout nœud est associé à une clé qui est :
supérieure à chaque clé de tous les nœuds de son \textbf{sous-arbre gauche} ;
inférieure à chaque clé de tous les nœuds de son \textbf{sous-arbre droit}.
\end{enumerate}
\end{reponse}
\begin{itemize}
\item (c) Dans un arbre binaire, un nœud peut avoir 0 enfant (c'est une feuille), 1 enfant ou 2 enfants. « Exactement deux » est donc faux : le nœud 5 de l'arbre 1 n'a qu'un enfant, et c'est bien un arbre binaire.
\item (d) Dans un ABR, les petites clés sont à gauche et les grandes à droite. La clé d'un nœud est donc \emph{plus grande} que tout ce qui est à sa gauche et \emph{plus petite} que tout ce qui est à sa droite.
\end{itemize}
\begin{erreur}
Inverser gauche et droite au point (d). Relis la phrase : « supérieure à chaque clé de son\ldots{} » : ce sont les clés \emph{plus petites} que le nœud, donc celles du sous-arbre gauche.
\end{erreur}

\Q{1 (b)}{3 points}{Pour chacun des deux arbres : (a) est-ce un arbre binaire ? (b) est-ce un ABR ? (c) sa hauteur ; (d) son parcours préfixe ; (e) son parcours infixe ; (f) son parcours suffixe.}
\begin{reponse}
\renewcommand{\arraystretch}{1.25}
\begin{tabular}{@{}p{3.2cm}p{6.2cm}p{5.6cm}@{}}
 & \textbf{Arbre 1} & \textbf{Arbre 2} \\
(a) arbre binaire ? & oui : chaque nœud a au plus 2 enfants & oui : chaque nœud a au plus 2 enfants \\
(b) ABR ? & non : 7 est à gauche de 6, alors que $7 > 6$ & oui \\
(c) hauteur & 3 & 2 \\
(d) préfixe & 13, 6, 7, 5, 1, 4 & 6, 4, 3, 5, 8, 7, 9 \\
(e) infixe & 7, 6, 5, 1, 13, 4 & 3, 4, 5, 6, 7, 8, 9 \\
(f) suffixe & 7, 1, 5, 6, 4, 13 & 3, 5, 4, 7, 9, 8, 6 \\
\end{tabular}
\end{reponse}

\Qp{(a) Arbre binaire.}{Un arbre est binaire si chaque nœud a au plus deux enfants.}
Compte les enfants de chaque nœud. Arbre 1 : 13 en a 2, 6 en a 2, 5 en a 1, les autres 0. Arbre 2 : 6, 4 et 8 en ont 2, les autres 0. Aucun nœud n'a 3 enfants ou plus : les deux arbres sont binaires.

\Qp{(b) Arbre binaire de recherche.}{Pour chaque nœud, tout son sous-arbre gauche doit être plus petit et tout son sous-arbre droit plus grand.}
\begin{itemize}
\item \textbf{Arbre 1.} Le nœud 6 a pour enfant gauche 7, et $7 > 6$ : la règle n'est pas respectée. Il y a d'autres problèmes : à droite de 13, on trouve 4, alors que $4 < 13$. Une seule erreur suffit pour répondre « non », mais il faut la citer.
\item \textbf{Arbre 2.} À gauche de 6 : 3, 4, 5 (tous $< 6$) ; à droite : 7, 8, 9 (tous $> 6$). À gauche de 4 : 3 ; à droite : 5. À gauche de 8 : 7 ; à droite : 9. Tout est respecté.
\end{itemize}
\textbf{Contrôle.} Le parcours infixe de l'arbre 2 est croissant (3, 4, 5, 6, 7, 8, 9) : c'est la signature d'un ABR. Celui de l'arbre 1 ne l'est pas.

\Qp{(c) Hauteur.}{On cherche le plus long chemin de la racine à une feuille, en comptant les arêtes.}
Arbre 1 : 13 $\to$ 6 $\to$ 5 $\to$ 1, soit 3 arêtes. Arbre 2 : 6 $\to$ 4 $\to$ 3, soit 2 arêtes. (Avec la convention qui compte les nœuds : 4 et 3.)

\Qp{(d), (e), (f) Parcours en profondeur.}{Les trois parcours appliquent la même règle à chaque nœud ; seul change le moment où l'on écrit le nœud.}
\begin{center}
\renewcommand{\arraystretch}{1.25}
\begin{tabular}{|l|l|}
\hline
\textbf{Préfixe} & nœud, puis sous-arbre gauche, puis sous-arbre droit \\ \hline
\textbf{Infixe} & sous-arbre gauche, puis nœud, puis sous-arbre droit \\ \hline
\textbf{Suffixe} & sous-arbre gauche, puis sous-arbre droit, puis nœud \\ \hline
\end{tabular}
\end{center}
\application Préfixe de l'arbre 1, pas à pas :
\begin{enumerate}
\item on écrit la racine \textbf{13}, puis on traite tout le sous-arbre gauche (racine 6), puis le sous-arbre droit (racine 4) ;
\item sous-arbre de 6 : \textbf{6}, puis son gauche \textbf{7}, puis son droit (racine 5) : \textbf{5}, puis \textbf{1} (5 n'a pas d'enfant gauche) ;
\item sous-arbre de 4 : \textbf{4}.
\end{enumerate}
On obtient 13, 6, 7, 5, 1, 4. Pour l'infixe, 5 n'a pas d'enfant gauche : on écrit 5 \emph{avant} son enfant droit 1, d'où « 5, 1 ».
\begin{erreur}
Dans l'infixe de l'arbre 1, écrire « 1, 5 » parce que 1 est en dessous de 5. Or 1 est l'enfant \emph{droit} de 5 : il passe après 5. La position à gauche ou à droite compte, pas la hauteur sur le dessin.
\end{erreur}

\begin{remarque}
Dans le PDF du sujet, cette sous-question est numérotée « (e) » au lieu de « (b) » : les listes imbriquées partagent le même compteur. C'est bien la deuxième partie de la question 1.
\end{remarque}

\section*{2. Programmation et implémentation des arbres}

\Q{2}{1 point}{Donner le code permettant d'implémenter l'arbre 2.}
\begin{reponse}
\begin{lstlisting}
n3 = Noeud(3)
n5 = Noeud(5)
n7 = Noeud(7)
n9 = Noeud(9)
n4 = Noeud(4, n3, n5)
n8 = Noeud(8, n7, n9)
arbre2 = Noeud(6, n4, n8)
\end{lstlisting}
\end{reponse}
\methode On imite le code donné pour l'arbre 1.
\begin{enumerate}
\item L'appel \code{Noeud(v, g, d)} crée un nœud de valeur \code{v}, d'enfant gauche \code{g} et d'enfant droit \code{d}. Sans \code{g} ni \code{d}, ils valent \code{None} : le nœud est une feuille.
\item On crée d'abord les \textbf{feuilles} (lignes 1 à 4), car un nœud a besoin que ses enfants existent déjà pour les recevoir.
\item On remonte ensuite étage par étage : les nœuds 4 et 8 (lignes 5 et 6), puis la racine 6 (ligne 7).
\item L'ordre des paramètres compte : \code{Noeud(4, n3, n5)} met 3 à gauche et 5 à droite.
\end{enumerate}
\textbf{Variante} acceptée, en une seule instruction :
\begin{lstlisting}
arbre2 = Noeud(6, Noeud(4, Noeud(3), Noeud(5)),
                  Noeud(8, Noeud(7), Noeud(9)))
\end{lstlisting}

\Q{3}{1 point}{(a) Résultat de l'expression \code{arbre1.gauche.droite.valeur} ? (b) Commande pour accéder à la valeur du nœud 1 depuis la racine.}
\begin{reponse}
(a) Telle qu'elle est écrite, l'expression provoque une erreur :
\code{AttributeError: 'Noeud' object has no attribute 'droite'}, car l'attribut s'appelle \code{droit}. \\
L'expression attendue, \code{arbre1.gauche.droit.valeur}, vaut \textbf{5}. \\
(b) \code{arbre1.gauche.droit.droit.valeur}
\end{reponse}
\methode On lit l'expression de gauche à droite, un point à la fois, en suivant le dessin.
\application
\begin{itemize}
\item \code{arbre1} est le nœud 13 ;
\item \code{arbre1.gauche} est son enfant gauche : le nœud 6 ;
\item \code{arbre1.gauche.droit} est l'enfant droit de 6 : le nœud 5 ;
\item \code{.valeur} donne sa valeur : 5.
\end{itemize}
Pour atteindre 1 : 13, à gauche 6, à droite 5, à droite 1. D'où l'expression suivante, qui vaut 1 :
\begin{center}\code{arbre1.gauche.droit.droit.valeur}\end{center}

\textbf{L'erreur affichée.} Le message \code{AttributeError} se traduit par « erreur d'attribut » : Python dit qu'un objet de la classe \code{Noeud} n'a pas d'attribut nommé \code{droite}. Les attributs sont exactement ceux créés dans \code{\_\_init\_\_} : \code{gauche}, \code{droit} et \code{valeur}. Un seul caractère de différence suffit pour provoquer l'erreur.
\begin{remarque}
L'énoncé écrit \code{droite} au lieu de \code{droit}. On attend sans doute la valeur 5 ; la réponse « le code provoque une \code{AttributeError} » est aussi juste, si elle est expliquée.
\end{remarque}

\Q{4}{1 point}{Compléter la fonction \code{est\_feuille} qui teste si un objet de type \code{Noeud} est une feuille.}
\begin{reponse}
\begin{lstlisting}
def est_feuille(abr):
    if abr.gauche is None and abr.droit is None:
        return True
    else:
        return False
\end{lstlisting}
\end{reponse}
\begin{itemize}
\item Ligne 2 : une feuille est un nœud sans enfant, donc sans enfant gauche \textbf{et} sans enfant droit. Dans la classe \code{Noeud}, « pas d'enfant » s'écrit \code{None}.
\item Ligne 3 : si les deux enfants sont absents, c'est une feuille : on renvoie \code{True}.
\item Ligne 5 : sinon, le nœud a au moins un enfant : on renvoie \code{False}.
\end{itemize}
Exemples avec l'arbre 1 : \code{est\_feuille(arbre1)} renvoie \code{False} ; \code{est\_feuille(arbre1.droit)} (le nœud 4) renvoie \code{True} ; \code{est\_feuille(arbre1.gauche.droit)} (le nœud 5, qui a un enfant droit) renvoie \code{False}.
\textbf{Variante} plus courte : \code{return abr.gauche is None and abr.droit is None}.
\begin{erreur}
Écrire \code{or} au lieu de \code{and}. Avec \code{or}, le nœud 5 (pas d'enfant gauche, mais un enfant droit) serait pris pour une feuille. Autre erreur : écrire \code{abr.gauche = None}, avec un seul \code{=} : c'est une affectation, pas une comparaison, et Python signale une erreur de syntaxe dans un \code{if}.
\end{erreur}

\section*{3. Manipulation récursive des arbres}

\Q{5}{3 points}{(a) La fonction \code{mystere} est-elle récursive ? (b) Valeur de retour des appels \code{mystere(abr1, 9)}, \code{mystere(abr1, 5)}, \code{mystere(abr2, 6)} et \code{mystere(abr2, 7)}. (c) Que fait la fonction ?}
\begin{lstlisting}
def mystere(abr, x):
    if abr is None :
        return 0
    else :
        return 1 + mystere(abr.gauche, x) + mystere(abr.droit, x)
\end{lstlisting}
\begin{reponse}
(a) \textbf{Oui} : à la ligne 5, la fonction \code{mystere} s'appelle elle-même, sur le sous-arbre gauche et sur le sous-arbre droit. \\
(b) \code{mystere(abr1, 9)} : \textbf{6} ; \code{mystere(abr1, 5)} : \textbf{6} ; \code{mystere(abr2, 6)} : \textbf{7} ; \code{mystere(abr2, 7)} : \textbf{7}. \\
(c) \code{mystere(abr, x)} renvoie le \textbf{nombre de nœuds} de l'arbre \code{abr}, c'est-à-dire sa \textbf{taille}. Le paramètre \code{x} ne sert à rien : il n'est jamais comparé à quoi que ce soit.
\end{reponse}
\textbf{(a) Récursivité.} Une fonction est \textbf{récursive} si elle s'appelle elle-même. Elle a aussi un \textbf{cas de base}, qui s'arrête sans appel récursif : ici, l'arbre vide (lignes 2 et 3), qui renvoie 0.

\textbf{(b) Le fonctionnement, bloc par bloc.}
\begin{itemize}
\item Lignes 2--3 : un arbre vide n'a aucun nœud : on renvoie 0.
\item Ligne 5 : sinon, on compte 1 pour la racine, plus le nombre de nœuds du sous-arbre gauche, plus celui du sous-arbre droit.
\end{itemize}
\textbf{Trace de} \code{mystere(abr1, 9)} (l'arbre 1, de racine 13). On note \code{m(n)} l'appel sur le nœud n ; \code{x} vaut 9 partout et ne change rien :
\begin{center}
\renewcommand{\arraystretch}{1.2}
\begin{tabular}{|l|l|c|}
\hline
\textbf{Appel} & \textbf{Calcul} & \textbf{Renvoie} \\ \hline
\code{m(7)} & $1 + 0 + 0$ (deux enfants vides) & 1 \\
\code{m(1)} & $1 + 0 + 0$ & 1 \\
\code{m(5)} & $1 + 0 + \code{m(1)} = 1 + 0 + 1$ & 2 \\
\code{m(6)} & $1 + \code{m(7)} + \code{m(5)} = 1 + 1 + 2$ & 4 \\
\code{m(4)} & $1 + 0 + 0$ & 1 \\
\code{m(13)} & $1 + \code{m(6)} + \code{m(4)} = 1 + 4 + 1$ & \textbf{6} \\ \hline
\end{tabular}
\end{center}
Python lance d'abord \code{m(13)}, qui a besoin de \code{m(6)}, qui a besoin de \code{m(7)} et de \code{m(5)}, et ainsi de suite. Les résultats remontent ensuite depuis les feuilles : c'est pourquoi le tableau commence par elles.

Comme \code{x} n'est pas utilisé, \code{mystere(abr1, 5)} donne aussi 6. L'arbre 2 a 7 nœuds (6, 4, 8, 3, 5, 7, 9) : \code{mystere(abr2, 6)} et \code{mystere(abr2, 7)} renvoient 7.

\textbf{(c)} Une bonne réponse dit ce que la fonction \textbf{renvoie} : « la taille de l'arbre (son nombre de nœuds) ». Elle peut ajouter que \code{x} est inutile.
\begin{erreur}
Croire que la fonction cherche \code{x}, parce qu'elle ressemble à une fonction de recherche, et répondre \code{True} ou \code{False}. La fonction renvoie des nombres (\code{return 0} et \code{return 1 + \ldots}) : regarde toujours ce qui est renvoyé avant de conclure. Autre erreur : compter les arêtes (5 pour l'arbre 1) au lieu des nœuds (6).
\end{erreur}
\begin{remarque}
Le paramètre \code{x} ne sert à rien dans cette fonction. C'est sans doute voulu, pour vérifier qu'on lit vraiment le code. Une version sans ce paramètre, \code{taille(abr)}, ferait exactement le même calcul.
\end{remarque}


\section*{4. Les rotations}

\begin{retenir}[Rotation droite, en trois étapes (code du sujet)]
\begin{enumerate}
\item \code{pivot = abr.gauche} : le pivot est l'enfant gauche de la racine.
\item \code{abr.gauche = pivot.droit} : le sous-arbre droit du pivot devient le sous-arbre gauche de la racine.
\item \code{pivot.droit = abr} : l'ancienne racine devient l'enfant droit du pivot. Le pivot est la nouvelle racine.
\end{enumerate}
La rotation gauche fait exactement la même chose en échangeant partout gauche et droit.
\end{retenir}

\Q{6 (a)}{1 point}{Dessiner l'arbre obtenu après une rotation droite de l'arbre 2.}
\begin{reponse}
\begin{center}
\begin{tikzpicture}[arbre, level 1/.style={sibling distance=26mm}, level 2/.style={sibling distance=14mm}, level 3/.style={sibling distance=10mm}]
\node {4} child {node {3}} child {node {6} child {node {5}} child {node {8} child {node {7}} child {node {9}}}};
\end{tikzpicture}
\end{center}
\end{reponse}
\application à l'arbre 2 (racine 6) :
\begin{enumerate}
\item le pivot est l'enfant gauche de 6 : le nœud \textbf{4} (avec ses enfants 3 et 5) ;
\item le sous-arbre droit du pivot (le nœud \textbf{5}) devient l'enfant gauche de 6 ;
\item le nœud 6 (avec maintenant 5 à gauche et toujours 8, 7, 9 à droite) devient l'enfant droit de 4. Le nœud 3 reste l'enfant gauche de 4.
\end{enumerate}
\textbf{Contrôle.} L'infixe est toujours 3, 4, 5, 6, 7, 8, 9 : la rotation a conservé la propriété d'ABR, comme annoncé dans le sujet.

\Q{6 (b)}{1 point}{Dessiner l'arbre obtenu après une rotation gauche de l'arbre 1.}
\begin{reponse}
\begin{center}
\begin{tikzpicture}[arbre, level 1/.style={sibling distance=26mm}, level 2/.style={sibling distance=20mm}, level 3/.style={sibling distance=13mm}]
\node {4} child {node {13} child {node {6} child {node {7}} child {node {5} child[missing] child {node {1}}}} child[missing]} child[missing];
\end{tikzpicture}
\end{center}
\end{reponse}
\application à l'arbre 1 (racine 13), avec l'algorithme de rotation gauche :
\begin{enumerate}
\item le pivot est le sous-arbre droit de la racine : le nœud \textbf{4} ;
\item le sous-arbre gauche du pivot devient le sous-arbre droit de 13. Ici, 4 n'a pas d'enfant gauche : 13 n'a plus d'enfant droit ;
\item le nœud 13 (avec tout son sous-arbre gauche 6, 7, 5, 1) devient l'enfant gauche de 4, qui est la nouvelle racine. Le nœud 4 n'a pas d'enfant droit.
\end{enumerate}
\textbf{Remarque.} L'arbre 1 n'est pas un ABR, donc le résultat n'en est pas un non plus. La question porte seulement sur la transformation : on applique les étapes sans chercher à ranger les clés. On vérifie qu'on n'a perdu aucun nœud : il y en a toujours 6.
\begin{erreur}
Oublier le sous-arbre gauche de 13 en route. Une rotation ne fait que déplacer deux liens : tout ce qui était sous 6 reste sous 6.
\end{erreur}

\Q{6 (c)}{2 points}{Écrire le code d'une fonction \code{rotation\_gauche} qui réalise la rotation gauche d'un ABR autour de sa racine.}
\begin{reponse}
\begin{lstlisting}
def rotation_gauche(abr):
    pivot = abr.droit
    abr.droit = pivot.gauche
    pivot.gauche = abr
    return Noeud(pivot.valeur, pivot.gauche, pivot.droit)
\end{lstlisting}
\end{reponse}
\begin{remarque}
L'énoncé parle d'une « méthode » ; ici, comme \code{rotation\_droite}, c'est une \emph{fonction} : elle n'est pas définie dans la classe \code{Noeud} et reçoit l'arbre en paramètre.
\end{remarque}
\methode On part de \code{rotation\_droite} et on échange \code{gauche} et \code{droit} partout. Chaque ligne suit une étape de l'algorithme du sujet :
\begin{itemize}
\item ligne 2 : le pivot est le sous-arbre droit de la racine ;
\item ligne 3 : le sous-arbre gauche du pivot devient le sous-arbre droit de la racine ;
\item ligne 4 : la racine ainsi modifiée devient le sous-arbre gauche du pivot ;
\item ligne 5 : on renvoie un nœud de valeur \code{pivot.valeur}, avec les mêmes enfants que le pivot : c'est la nouvelle racine.
\end{itemize}
\textbf{Exemple.} Avec l'ABR de racine 20, d'enfants 15 et 30, et 25, 40 sous 30 : \code{rotation\_gauche} renvoie l'arbre de racine 30, de préfixe 30, 20, 15, 25, 40. Son infixe 15, 20, 25, 30, 40 est toujours croissant.

\textbf{Variante.} La ligne 5 peut être remplacée par \code{return pivot} : le pivot est déjà la nouvelle racine, avec les bons enfants.
\begin{erreur}
Inverser les lignes 3 et 4. Si on écrit d'abord \code{pivot.gauche = abr}, l'ancien sous-arbre gauche du pivot est perdu : la ligne suivante ne peut plus le récupérer. Autre oubli : le \code{return} final. Sans lui, la fonction renvoie \code{None}.
\end{erreur}

\section*{5. Graphes}

\begin{retenir}[Le vocabulaire des graphes]
Un \textbf{graphe} est fait de \textbf{sommets} (ici, les villes) reliés par des \textbf{arêtes} (ici, les vols directs). Le graphe est \textbf{non orienté} : un vol de Paris à Berlin existe aussi de Berlin à Paris. Il est \textbf{pondéré} : chaque arête porte un nombre, son \textbf{poids} (ici, la distance en milliers de kilomètres).
\end{retenir}

\Q{7}{1 point}{Représenter le graphe des vols de la compagnie AirInfo.}
\begin{reponse}
\begin{center}
\begin{tikzpicture}[every node/.style={draw,circle,minimum size=8mm,font=\bfseries}, thick, x=1cm, y=1cm]
\node (W) at (0,0) {W};
\node (P) at (3,1.6) {P};
\node (B) at (3,-1.6) {B};
\node (T) at (7,-1.6) {T};
\node (S) at (7,1.6) {S};
\begin{scope}[every node/.style={fill=white, inner sep=1.5pt, font=\small}]
\draw (W) -- node {9} (P);
\draw (W) -- node {4} (B);
\draw (P) -- node {3} (B);
\draw (P) -- node[pos=0.35] {12} (T);
\draw (P) -- node {17} (S);
\draw (B) -- node {7} (T);
\draw (T) -- node {7} (S);
\end{scope}
\end{tikzpicture}
\end{center}
\end{reponse}
\methode
\begin{enumerate}
\item Dessine un sommet (un cercle) par ville : W, P, B, T, S.
\item Pour chaque ligne de la liste des distances, trace une arête entre les deux villes et écris la distance dessus.
\item Vérifie à la fin : la liste compte 7 lignes, ton graphe doit avoir 7 arêtes.
\end{enumerate}
La place des sommets sur la feuille n'a pas d'importance : seules comptent les arêtes et leurs poids. Essaie de placer les sommets pour que les arêtes se croisent le moins possible.
\begin{remarque}
L'énoncé dit que la compagnie propose « certains de ces vols » sans préciser lesquels. Le code de la question 8 (\code{'W': \{'P': 9, 'B': 4\}}, \code{\{'P': 17, 'T': 7\}}) montre qu'on garde les 7 vols de la liste : c'est ce graphe qui est attendu.
\end{remarque}

\Q{8}{1 point}{Compléter le code du dictionnaire \code{graphe\_airinfo}.}
\begin{reponse}
\begin{lstlisting}
graphe_airinfo = {
    'W': {'P': 9, 'B': 4},
    'P': {'W': 9, 'B': 3, 'T': 12, 'S': 17},
    'B': {'W': 4, 'P': 3, 'T': 7},
    'T': {'P': 12, 'B': 7, 'S': 7},
    'S': {'P': 17, 'T': 7}
}
\end{lstlisting}
\end{reponse}
\methode
\begin{enumerate}
\item Chaque clé du grand dictionnaire est un sommet : \code{'W'}, \code{'P'}, \code{'B'}, \code{'T'}, \code{'S'}.
\item Sa valeur est un petit dictionnaire : ses clés sont les voisins du sommet, ses valeurs les distances.
\item Le graphe est non orienté : chaque arête apparaît \textbf{deux fois}. Le vol Paris -- Berlin de 3 donne \code{'B': 3} dans la ligne de P \textbf{et} \code{'P': 3} dans la ligne de B.
\end{enumerate}
La dernière ligne, \code{\ldots: \{'P': 17, 'T': 7\}}, est celle du seul sommet relié à Paris par 17 et à Tokyo par 7 : Sydney, \code{'S'}.
\begin{erreur}
Oublier le « sens retour » d'une arête : par exemple écrire \code{'W'} dans la ligne de P mais pas \code{'P'} dans celle de W. Contrôle : la somme des longueurs des petits dictionnaires vaut $2 + 4 + 3 + 3 + 2 = 14$, soit deux fois 7 arêtes.
\end{erreur}

\Q{9}{1 point}{Donner la matrice d'adjacence du graphe.}
\begin{reponse}
Avec les sommets dans l'ordre W, P, B, T, S (lignes et colonnes) :
\[
\begin{array}{c|ccccc}
 & W & P & B & T & S \\ \hline
W & 0 & 1 & 1 & 0 & 0 \\
P & 1 & 0 & 1 & 1 & 1 \\
B & 1 & 1 & 0 & 1 & 0 \\
T & 0 & 1 & 1 & 0 & 1 \\
S & 0 & 1 & 0 & 1 & 0
\end{array}
\]
\end{reponse}
\methode
\begin{enumerate}
\item Choisis un ordre des sommets et écris-le au-dessus des colonnes et devant les lignes. Annonce-le dans ta réponse.
\item Case (ligne $i$, colonne $j$) : 1 s'il y a une arête entre $i$ et $j$, 0 sinon.
\item La diagonale est faite de 0 : aucune ville n'a de vol vers elle-même.
\end{enumerate}
\textbf{Contrôle.} Le graphe est non orienté, donc la matrice est \textbf{symétrique} : la case (P, B) est égale à la case (B, P). La somme d'une ligne donne le nombre de voisins : 2 pour W, 4 pour P.

\textbf{Variante} acceptée : la matrice \textbf{pondérée}, qui met la distance au lieu du 1 (ligne de W : 0, 9, 4, 0, 0 ; ligne de P : 9, 0, 3, 12, 17 ; ligne de B : 4, 3, 0, 7, 0 ; ligne de T : 0, 12, 7, 0, 7 ; ligne de S : 0, 17, 0, 7, 0).

\Q{10}{1 point}{Déterminer la valeur de l'expression \code{graphe\_airinfo['W']['P']}.}
\begin{reponse}
\textbf{9}
\end{reponse}
On lit l'expression de gauche à droite :
\begin{enumerate}
\item \code{graphe\_airinfo['W']} est la valeur associée à la clé \code{'W'} : le dictionnaire \code{\{'P': 9, 'B': 4\}} ;
\item \code{['P']} prend dans ce dictionnaire la valeur associée à \code{'P'} : 9.
\end{enumerate}
C'est la distance du vol Washington -- Paris : 9 000 km.

\Q{11}{1 point}{Écrire une fonction \code{vol\_direct(graphe, ville1, ville2)} qui renvoie \code{True} si les deux villes sont reliées par une arête, \code{False} sinon.}
\begin{reponse}
\begin{lstlisting}
def vol_direct(graphe, ville1, ville2):
    return ville2 in graphe[ville1]
\end{lstlisting}
\end{reponse}
\textbf{Explication.} \code{graphe[ville1]} est le dictionnaire des voisins de \code{ville1}. Le mot-clé \code{in} teste si \code{ville2} est une \textbf{clé} de ce dictionnaire, c'est-à-dire un voisin. L'expression vaut déjà \code{True} ou \code{False} : on la renvoie directement.

\textbf{Exemples.} \code{vol\_direct(graphe\_airinfo, 'T', 'B')} renvoie \code{True} (le vol Berlin -- Tokyo existe) ; \code{vol\_direct(graphe\_airinfo, 'W', 'T')} renvoie \code{False}.

\textbf{Variante} avec une boucle, plus longue mais aussi juste :
\begin{lstlisting}
def vol_direct(graphe, ville1, ville2):
    for voisin in graphe[ville1]:
        if voisin == ville2:
            return True
    return False
\end{lstlisting}
\begin{erreur}
Placer \code{return False} dans la boucle, dans un \code{else} : la fonction s'arrête alors dès le premier voisin, même si \code{ville2} est plus loin dans le dictionnaire. Le \code{return False} doit venir \textbf{après} la boucle, quand tous les voisins ont été examinés.
\end{erreur}
\begin{remarque}
Le deuxième exemple du sujet, \code{vol\_direct(graphe\_airinfo, 'W', 'B')}, annonce \code{False}. Or le vol Washington -- Berlin (4) est dans la liste et dans le code donné (\code{'W': \{'P': 9, 'B': 4\}}) : la fonction correcte renvoie \code{True}. Un exemple qui renvoie \code{False} est \code{vol\_direct(graphe\_airinfo, 'W', 'T')}.
\end{remarque}


\end{document}
